% ============================================================================
% SECTION 5 — RESULTS
% Paper budget: ~3 pages + 2 tables + 1 figure. Awaits the pilot and the
% sweep; the structure is fixed now so tables land in place.
% NOTE: subsections follow the CANONICAL RQ numbering from
% Section~\ref{sec:intro-rqs} (RQ1 ordering, RQ2 orchestration, RQ3
% open-weight). The old thesis skeleton used a superseded numbering.
% Data source when writing: results.jsonl via src/tdbf/metrics.py
% (build_report — the same numbers as script 08 and the webapp /metrics page).
% Discipline for this section: report, do not interpret. Every "because" and
% every "this suggests" belongs in Section~\ref{sec:discussion}.
% ============================================================================
\section{Results}
\label{sec:results}

\subsection{RQ1: Test-First versus Reactive}
\label{sec:results-rq1}
% TODO after the sweep. Lead with the headline number in the first sentence.
% Table I: resolved-rate per arm (n, resolved, rate, Δ, 95% bootstrap CI,
%          McNemar p) — overall, with a per-repository breakdown underneath.
% Figure:  per-instance paired outcome grid (A✓B✓ / A✓B✗ / A✗B✓ / A✗B✗) —
%          this figure is what makes the design visibly paired.
% Table II: cost and efficiency per arm (tokens, iterations, $/instance,
%          $/resolved).

\subsection{Test Generation Quality}
\label{sec:results-brt}
% TODO. L0–L4 distribution, C2P, and gate-accept rate over the sweep sample.
% Development-set numbers (C2P 5/6 = 0.83 on astropy) are labelled pilot and
% are not pooled with the sweep.

\subsection{RQ2: Orchestration Ablations}
\label{sec:results-rq2}
% TODO. One table, one row per variant, columns = resolved-rate, Δ vs. the
% default pipeline, and $/instance. Include the NEGATIVE results as rows, not
% as prose apologies:
%   - self-review (critic/reviewer) on a strong model: revise-rate 1/14, a
%     near no-op;
%   - post-fix adversarial verification: n=2, both false-negative "robust" —
%     depth-1 self-review without an oracle fails open;
%   - best-of-N test selection; attempt budgets.
% Also report here: the gold-free (L3∪L4) acceptance criterion vs. an
% oracle L4 gate — the ablation that justifies the design in
% Section~\ref{sec:method-phase1}.
% Offline artifact analysis (test properties vs. resolved) goes at the end of
% this subsection, explicitly labelled exploratory.

\subsection{RQ3: Open-Weight Models}
\label{sec:results-rq3}
% TODO. One table: per model — terminal-submission rate, reproduction rate,
% resolved-rate — so the breakdown stage is readable directly from the row.
% Known finding to report plainly: local 7–14B models call tools correctly but
% do not sustain the autonomous loop to a terminal submit; at 14B and a larger
% iteration budget they submit but fabricate expected values. Prompt-hardening
% and a larger local model did not close the gap.
